\subsection{Baseline reproduction: KAN-ODE vs.\ MLP-ODE}
\label{sec:kan-vs-mlp}

Two MLP vector fields with 252 parameters were trained under the identical
recipe: the paper's literal baseline, one hidden layer of 50 tanh units, and
a deeper $[2,14,8,8,2]$ SiLU network that trains reliably under our
protocol. Table~\ref{tab:kan-vs-mlp} and Figures~\ref{fig:kan-vs-mlp}
and~\ref{fig:kan-vs-mlp-t28} compare
them with the KAN-ODE. A fourth column, added after auditing the paper's own
code (Section~\ref{sec:audit}), retrains the identical literal baseline under
the paper's own learning rate and initialization scale instead of our
default. Figures~\ref{fig:kan-vs-mlp} and~\ref{fig:kan-vs-mlp-t28} plot that run
in place of the default-setting tanh MLP.

\begin{table*}[t]
\centering
\footnotesize
\renewcommand{\arraystretch}{1.12}
\caption{KAN-ODE against parameter-matched MLP-ODEs at $10^4$ epochs.
Epoch counts are the first epoch at which the training loss falls below the
threshold. The last three rows are computed from the saved checkpoints:
field error is $\|\mathbf{f}_\theta-\mathbf{f}\|_2$ averaged over the true
orbit on $(3.5,14]$ and divided by the mean $\|\mathbf{f}\|_2=4.82$ there;
the period is measured from zero-crossings of the predicted prey signal over
$[0,28]$; $H$ drift is $\max_t|H(\hat{\mathbf{u}}(t))-H_0|/H_0$ on $[0,28]$.
The fourth column is the identical literal baseline of the third, retrained
under the paper's own learning rate and initialization scale
(Section~\ref{sec:audit}); its $5\times10^4$-epoch continuation is in
Table~\ref{tab:phase4}.}
\label{tab:kan-vs-mlp}
\begin{tabularx}{\textwidth}{@{}Y C{2.85cm} C{2.85cm} C{2.85cm} C{2.85cm}@{}}
\toprule
\hdrrow \thh{Metric} & \thh{KAN-ODE (RBF)} & \thh{MLP-ODE (SiLU)} & \thh{MLP-ODE (tanh, default lr/init)} & \thh{MLP-ODE (tanh, paper's lr/init)}\\
\midrule
Parameters & 240 & 252 & 252 & 252\\
Training MSE & $8.83\times10^{-5}$ & $9.50\times10^{-5}$ & $1.08\times10^{0}$ & $1.11\times10^{-3}$\\
Extrapolation MSE & $\mathbf{8.92\times10^{-5}}$ & $4.77\times10^{-4}$ & $1.69\times10^{0}$ & $9.65\times10^{-3}$\\
Extrapolation $R^2$ & $\mathbf{0.99997}$ & $0.99984$ & $0.4215$ & $0.99669$\\
Relative $L_2$ error & $\mathbf{0.33\%}$ & $0.76\%$ & $45.2\%$ & $3.42\%$\\
Lipschitz estimate $\hat L$ & $6.91$ & $\mathbf{5.07}$ & $38.22$ & $9.23$\\
Epochs to $10^{-2}$ / $10^{-3}$ / $10^{-4}$ & $3112$ / $5282$ / $9617$ & $\mathbf{741}$ / $\mathbf{1297}$ / $\mathbf{9145}$ & never & $3461$ / never / never\\
Wall-clock, $10^4$ epochs & $1705$ s & $2856$ s & $736$ s & $1052$ s\\
\midrule
Field error on the orbit, $(3.5,14]$ & $2.84\%$ & $2.82\%$ & -- & $7.03\%$\\
Relative period error & $\mathbf{1.7\times10^{-5}}$ & $8.0\times10^{-4}$ & -- & $3.1\times10^{-3}$\\
Max.\ first-integral drift & $3.1\%$ & $\mathbf{2.8\%}$ & -- & $22.7\%$\\
\bottomrule
\end{tabularx}
\end{table*}

\begin{figure}[t]
\centering
\includegraphics[width=\columnwidth]{figures/ablation/03_kan_vs_mlp_convergence.png}
\caption{Training MSE per epoch for KAN-ODE and the paper's
parameter-matched tanh MLP-ODE, trained with the paper's own learning rate
and initialization (fourth column of Table~\ref{tab:kan-vs-mlp}). KAN-ODE
descends steadily and reaches $10^{-4}$ at epoch 9617; the tanh MLP reaches
$10^{-2}$ at epoch 3461 but never $10^{-3}$ within $10^4$ epochs.}
\label{fig:kan-vs-mlp}
\end{figure}

\begin{figure}[t]
\centering
\includegraphics[width=\columnwidth]{figures/ablation/07_extrapolation_t28.png}
\caption{Predicted against true populations out to $t=28$ for the three
models, with no retraining. KAN-ODE and SiLU-MLP are visually
indistinguishable from the truth. The paper's tanh MLP under its own
learning rate and initialization (bottom, fourth column of
Table~\ref{tab:kan-vs-mlp}) follows the training window and the first
periods closely, then runs slightly slow and lags visibly by the far
window $(14,28]$.}
\label{fig:kan-vs-mlp-t28}
\end{figure}

The fourth column shows that the paper's literal architecture was never the
problem. Given the paper's own learning rate and initialization scale, the
identical tanh network reaches $R^2=0.997$. We return to this in
Section~\ref{sec:phase4}, which also reports its $5\times10^4$-epoch
continuation. The next two paragraphs compare KAN-ODE against the working
SiLU baseline, columns one and two.

We find one result that confirms the paper and one that contradicts it.
\textbf{KAN-ODE extrapolates $5.4\times$ better} from a training loss only
$7\%$ lower, reproducing the paper's accuracy claim. The checkpoints show
where that gap comes from. Both fields are equally accurate on the orbit
($2.84\%$ against $2.82\%$ mean field error), and both hold the invariant
to about $3\%$, but the MLP's orbit period is off by $8.0\times10^{-4}$
against the KAN's $1.7\times10^{-5}$, a $48\times$ larger period error. On a
periodic orbit a period error $\epsilon_T$ becomes a phase lag that grows
linearly, $\approx 2\pi\epsilon_T\, t/T$, so the trajectory error of the MLP
keeps accumulating while the KAN's stays near its training-window level.
The extrapolation gap is a phase-accuracy gap rather than a field-accuracy
gap. \textbf{The MLP converges faster, not slower}: it reaches $10^{-3}$
training MSE at epoch 1297 against the KAN's 5282, $4.1\times$ sooner. This
contradicts the paper's headline efficiency claim on this same system (a
$10^4$-epoch KAN-ODE against a $10^5$-epoch MLP-ODE). Under our protocol, at
matched parameter count and optimizer settings, the MLP was the faster model
to train; KAN-ODE won on extrapolation, the paper's other headline claim.

\paragraph{Why the paper's literal MLP fails to train under our default
optimizer settings.} The $[2,50,2]$ tanh network never gets below a
training MSE of $1.08$ at our default learning rate
($2\times10^{-3}$) and initialization (Glorot, gain $1.0$). A controlled
2000-epoch factorial run, holding those optimizer settings fixed, separates
the two things that differ between it and the working baseline
(Table~\ref{tab:mlp-breakdown}).

\begin{table*}[t]
\centering
\footnotesize
\renewcommand{\arraystretch}{1.15}
\caption{Controlled breakdown of why the paper's literal MLP baseline
architecture, $[2,50,2]$ with hyperbolic tangent (tanh) activation, fails to
train at our default learning rate and initialization scale. All
four runs use 2000 epochs and Glorot initialization at gain $1.0$; none
varies the initialization scale (see the paragraph following this table).}
\label{tab:mlp-breakdown}
\begin{tabularx}{\textwidth}{@{}C{0.9cm} Y C{1.8cm} C{2.0cm} C{2.6cm} C{2.6cm} C{2.0cm}@{}}
\toprule
\hdrrow \thh{Run} & \thh{Architecture} & \thh{Activation} & \thh{Learning rate} & \thh{Training MSE} & \thh{Extrapolation MSE} & \thh{$R^2$}\\
\midrule
A & $[2,50,2]$ & tanh & $2\times10^{-3}$ & $1.21\times10^{0}$ & $1.98\times10^{0}$ & $0.319$\\
B & $[2,50,2]$ & tanh & $5\times10^{-4}$ & $1.85\times10^{0}$ & $9.23\times10^{0}$ & $-2.169$\\
C & $[2,14,8,8,2]$ & SiLU & $2\times10^{-3}$ & $\mathbf{4.28\times10^{-4}}$ & $\mathbf{6.97\times10^{-4}}$ & $\mathbf{0.99976}$\\
D & $[2,50,2]$ & SiLU & $2\times10^{-3}$ & $3.62\times10^{-3}$ & $2.68\times10^{2}$ & $-91.06$\\
\bottomrule
\end{tabularx}
\end{table*}

Runs A and D share the architecture and differ only in activation. The
training MSE falls from $1.21$ to $3.62\times10^{-3}$, a $335\times$
improvement, so at these optimizer settings, activation decides whether this
architecture fits the data at all. Runs C and D share the SiLU activation and
differ only in depth. The shallow network fits the window but its
extrapolation error explodes to $2.68\times10^{2}$, against
$6.97\times10^{-4}$ for the deep one, so depth decides whether a model that
trains also extrapolates. Lowering the learning rate alone (run B) makes the
tanh network worse, not better. Five times the epoch budget under these
same settings does not rescue it either: at $5\times10^4$ epochs, training
MSE moves only from $1.08$ to $1.01$ and extrapolation gets \emph{worse}
(Table~\ref{tab:phase4}).

\paragraph{The actual cause: learning rate and initialization scale, not
architecture.} None of runs A--D varies the initialization scale, and the
base paper's own published code turned out to use a combination this
factorial never tested: a learning rate of $10^{-2}$ (five times our
default) together with weights initialized near zero (Glorot
divided by $10^5$), for the identical $[2,50,2]$ tanh architecture. Retrained
with only those two settings changed, the literal architecture trains
(Table~\ref{tab:kan-vs-mlp}, fourth column). At $10^4$ epochs it reaches a
training MSE of $1.11\times10^{-3}$ (against $1.08$ before) and an
extrapolation MSE of $9.65\times10^{-3}$, $R^2=0.997$ (against $R^2=0.42$
before), with a Lipschitz estimate of $9.23$ against $38.2$. At $5\times10^4$
epochs it reaches $5.43\times10^{-5}$ training MSE
and $R^2=0.9998$ on extrapolation, with zero non-finite gradient steps and no
clipping in either run. Section~\ref{sec:phase4} reports the full comparison
against KAN-ODE and the SiLU MLP at both budgets, and updates the
implementation audit accordingly. The paper's literal recipe was never
architecturally incompatible with this system. Our default
optimizer settings, tuned around the KAN-ODE side of the comparison, were the
wrong settings for a wide, near-linear-at-initialization tanh network, and
Table~\ref{tab:mlp-breakdown}'s factorial measured sensitivity to activation
\emph{at those settings} rather than an architecture-level limit.
